%% ma: L11-B16-0002
%% lop: 11
%% chuong: 5
%% bai: 16
%% dang: 11.16.1
%% loai: TL
%% muc: VD
%% dap_an: "-"
%% nguon: "tai-lieu-hoc-tap-mon-toan-lop-11-theo-chuong-trinh-giao-duc-pho-thong-2018.pdf (Lớp Toán Thầy Duy, Phạm Lê Duy, THPT Chu Văn An), Chương 3 Bài 2, Dạng 1, Ví dụ 2 (tr.351)"
%% kien_thuc: "Lớp 11 Bài 16 (giới hạn của hàm số)"
%% trang_thai: cho-duyet
%% ly_do: "Dạng mới đề xuất (11.16.x), chờ giáo viên duyệt dạng; tài liệu gốc không có lời giải mẫu, đáp án do kho tự giải và kiểm bằng sympy."
\begin{ex}
Tìm giới hạn

a) $\lim\limits_{x\to-3}\dfrac{x+3}{x^2-9}$;

b) $\lim\limits_{x\to2}\dfrac{x^2+x-6}{2x^2-5x+2}$;

c) $\lim\limits_{x\to4}\dfrac{2x^2-7x-4}{x^2+x-20}$;

d) $\lim\limits_{x\to2}\dfrac{x^3+3x^2-9x-2}{-x^3+x+6}$;

e) $\lim\limits_{x\to-1}\dfrac{4x^2-3x-7}{x^3+1}$;

f) $\lim\limits_{x\to2}\dfrac{(x^2+x-6)^2}{-x^3+2x^2}$.
\loigiai{\sloppy a) $\dfrac{x+3}{(x-3)(x+3)}=\dfrac1{x-3}\to-\dfrac16$.

b) $\dfrac{(x-2)(x+3)}{(2x-1)(x-2)}=\dfrac{x+3}{2x-1}\to\dfrac53$.

c) $\dfrac{(2x+1)(x-4)}{(x+5)(x-4)}=\dfrac{2x+1}{x+5}\to1$.

d) $x^3+3x^2-9x-2=(x-2)(x^2+5x+1)$ và $-x^3+x+6=-(x-2)(x^2+2x+3)$, nên giới hạn $=-\dfrac{4+10+1}{4+4+3}=-\dfrac{15}{11}$.

e) $4x^2-3x-7=(x+1)(4x-7)$ và $x^3+1=(x+1)(x^2-x+1)$, nên giới hạn $=\dfrac{-4-7}{1+1+1}=-\dfrac{11}3$.

f) $x^2+x-6=(x-2)(x+3)$ và $-x^3+2x^2=-x^2(x-2)$, nên biểu thức $=\dfrac{(x-2)(x+3)^2}{-x^2}\to0$.}
\end{ex}
